Nguồn: \href{http://vanviet.info/van-de-hom-nay/nho-mot-phan-dnh-dieu-cuong-truc/}{Văn Việt} và \href{http://www.hasiphu.com/baivietmoi_203.html}{Thư viện Hà Sĩ Phu} 

Từ Tiệp Khắc về tôi được phân công vào Viện Sinh học thuộc Viện Khoa học Việt Nam mà GS Phan Đình Diệu là Viện phó, tương đương cấp Thứ trưởng. Tuy cùng cơ quan lớn nhưng tôi và GS Diệu chỉ ngồi chuyện trò riêng với nhau được một lần, tôi nhớ khoảng sau năm 2000. Lúc ấy tôi đã bị liệt vào tội nặng (khởi tố tội phản quốc nhưng không đem xử), còn bác Diệu thì từ lâu đã nổi tiếng là một đại biểu Trí thức trong Đoàn Chủ tịch Mặt trận Tổ quốc Việt Nam dám có những quan điểm mới lạ, không thuận chiều CS lúc bấy giờ, và ít nhiều cũng đã được ‘hưởng’ sự phê phán, nhất là sau cuộc trả lời phỏng vấn của nhà sử học Na Uy, ông Stein Tønnesson. GS Diệu nhìn sự phi lý của Mác-Lê dưới nhãn quan Toán học cũng như tôi nhìn sự phi lý ấy dưới nhãn quan Sinh học. Chúng tôi có sự đồng cảm tự nhiên.

Khi trò chuyện tôi truyền đạt ý kiến của một số anh em trí thức:

- Xã hội cần anh làm một Sakharov, anh ạ!

Bác Diệu cười:

- Mình chỉ biết làm Toán, tư duy Toán, từ đó cũng soi rọi vào xã hội, chứ không thích làm quan, làm chính trị hay làm một vai trò gì. Vai trò như vậy để Hà Sĩ Phu (sic).

Tôi lại càng cười:

- Anh nói đùa vậy thôi, tôi càng là anh bạch diện thư sinh dốt về chính trị, sức mọn tài hèn, biết gì mà làm?

Có lẽ do sự đồng cảm ấy, nên có lần mở ra cuộc thảo luận “Trí thức là ai” năm 2008, sau ý kiến của GS Diệu tôi đã viết một bài hưởng ứng (*), nói rõ mỗi người có thể nghiên cứu một lĩnh vực chuyên môn hẹp nhưng “chất Trí thức” thì không bao giờ bị giới hạn trong lĩnh vực hẹp ấy mà luôn xâm nhập vào mọi mặt của đời sống, nhất là những vấn nạn cấp thiết của đất nước mình, của nhân dân mình. Người Trí thức tham gia vào Chính trị trước hết bằng cách bộc lộ những nhận thức khoa học, khúc triết của mình.

Nay GS Phan Đình Diệu đã từ biệt chúng ta, về cõi vĩnh hằng, tôi xin thay mặt một số Trí thức ở Đà Lạt, nhóm Thân hữu Đà Lạt và CLB Phan Tây Hồ, có một Câu đối viếng như một nén nhang muộn kính cẩn thành tâm tiễn biệt một người Trí thức hiếm có của đất nước:

Trí tuệ hiếm hoi thay, ánh Toán học sáng bừng đôi mắt bác!

Tâm hồn cương trực thế, giữa Quan trường trơ trọi một mình ông!

H.S.P. (15-5-2018)

——————————————————————————————–

(*) Đây là bài viết hưởng ứng ý kiến của GS Phan Đình Diệu về giới Trí thức.

Bài viết như sau:


Công-Nông-Trí và Nguyên khí quốc gia

(Tham gia cuộc thảo luận: Trí thức là ai?)

Hà Sĩ Phu

(\url{http://www.hasiphu.com/baivietmoi_25.html}).

Ở Việt Nam ta, dưới sự lãnh đạo của Đảng Cộng sản, năm 1930 Trí thức từng được xếp số 1 trong chuỗi bốn kẻ thù TRÍ PHÚ ĐỊA HÀO cần “đào tận gốc, trốc tận rễ”. Nhưng rồi thế cục xoay vần, năm 2008 này, sau một khoá họp trung ương của Đảng, Trí thức chính thức được tái xác nhận rằng đã được kéo lên, xếp hàng cuối trong chuỗi ưu tú “CÔNG NÔNG TRÍ”. Chừng ấy năm trời lên được một bậc kể cũng khiêm nhường! Nhưng từ “kẻ thù” số 1 chuyển lên “đồng chí” số 3 cũng quý lắm chứ?

Lịch sử đã trôi 78 năm, tức đã quá một đời người được gọi “xưa nay hiếm”, hôm nay chúng ta đang thảo luận để tìm một định nghĩa cho loại công dân “quái dị” này. Không “quái dị” sao được, khi loại công dân này luôn làm cho Đảng đau đầu, trước đây đã bị “xuống chó” nay lại được “lên voi”, dẫu loại “voi” nhỏ con này vẫn rất cần có quản tượng cầm búa cầm liềm đứng bên chế ngự (không cần đọc “nghị quyết” tôi cũng biết chắc sự thể sẽ vẫn như thế).

Để có định nghĩa đúng cho một danh từ, một mặt cần hiểu từ nguyên, hiểu được xuất xứ. Nhưng mặt khác, từ ngữ là một thể sống chứ không chết cứng, bởi từ ngữ chỉ là cái vỏ quy ước mà cái hồn bên trong chính là nội dung khái niệm mà từ ngữ ấy bao hàm. Khái niệm thì luôn được làm giàu, được tu chỉnh theo sự phát triển của nhận thức, theo sự phát triển của xã hội. Vậy định nghĩa của một từ ngữ cũng phát triển theo thời gian, điều này không có gì lạ.

Dịp này, để tiến tới tìm một định nghĩa so sánh chính xác cho các thuật ngữ gốc Hán (trí thức, thức giả, học giả, hiền tài, sỹ phu…), gốc phương Tây (intellectuel, savant…), gốc Việt (kẻ Sĩ, kẻ có học, người lao động trí óc…) tôi mạn phép nêu lên vài nhận xét chủ quan, chắc hẳn còn nông cạn, về đặc điểm của loại hình công dân đặc biệt này, mà ta tạm dùng chữ Trí thức làm đại diện.

Trí thức có ba đặc điểm, hay ba tính chất:

1/ Tính chất tự do và không ranh giới: Người trí thức có thể đi sâu vào một lĩnh vực nhất định, nhưng không bị khoanh vùng trong lĩnh vực đó, trong quốc gia đó. Sự phân chia thành các lĩnh vực khoa học tự nhiên, khoa học xã hội, văn hoá, văn nghệ, chính trị, đạo đức, kinh tế, quốc gia, quốc tế, con người, nhân loại… vân vân… đều chỉ là những sự phân chia tương đối, quy ước, tạm thời. Trước sự phát triển và liên tưởng của lô-gích tư duy những ranh giới ấy đều bị phá bỏ.

2/ Tính chất tiên phong và phát hiện: Bởi được trang bị bằng hai công cụ đặc biệt là kho tri thức ngàn đời của nhân loại đã đúc kết thành các khoa học và sức nhạy cảm chủ quan của bộ óc cá nhân, những người Trí thức như những bộ máy dò tiên phong vô cùng nhạy cảm, nó bay về muôn phương, nó cảm ứng lập tức với những cái mới, cái bất thường, nó xâu chuỗi lập tức những quan hệ tuyến tính, những quan hệ thuận quy luật và trái quy luật. Học vấn và lao động trí óc tuy là những điều kiện cần thiết tất yếu cho đặc điểm này, nhưng cũng chỉ là những điều kiện làm nền, điều kiện ban đầu thôi, chưa hề đầy đủ để tạo nên Trí thức.

3/ Tính lương thiện, tính nhân văn hay tính “Người” (“người khôn ngoan” Homo sapiens): Con người là sự diễn tiến liên tục và bất tận từ con vật hoang dã hướng đến con người hoàn thiện. Trong đó, Trí thức là bộ máy lọc, chuyên lọc bỏ những cặn bã phi nhân tính. Sự thôi thúc của Nhân tính này khiến Trí thức cứ vơ nỗi đau của người khác thành nỗi đau của mình, mà tự làm khổ mình, tự gây rắc rối cho mình. Tất nhiên bộ máy lọc này cũng đang tiến hoá theo thời gian nên sự lọc này không thể thoát khỏi sự hạn chế của giai đoạn lịch sử, có thể lọc rồi sẽ lại được lọc nữa.

Tầm Trí thức của một người là do ba đặc điểm ấy có được ở cấp độ nào. Trí thức là chiếc máy dò vạn năng, vượt mọi ranh giới, là bộ máy lọc lương thiện rất mẫn cảm, và khi bắt được sự cảm ứng, sự cộng hưởng, thì những cảm ứng ấy lập tức biến thành năng lượng của lý trí hoặc của cảm xúc, thành những kết quả lý thuyết và thành hành động. Điều ấy giải thích tại sao những người thầy giỏi, có tư duy minh triết luôn đi kèm với nhân cách, với khí phách và sự dấn thân.

Ngoài sự thôi thúc của tính thiện, sự thôi thúc của lý tính tìm kiếm vô giới hạn cũng khiến Trí thức thường “xớ rớ” vào những việc không phải của mình, nhưng Trí thức muôn đời cứ “vô duyên” như thế!

Ba tiêu chuẩn “Nhân, Trí, Dũng” (hay Bi, Trí, Dũng) là ba tiêu chuẩn lý tưởng của con người ưu tú, cũng là tiêu chuẩn lý tưởng cho tầng lớp tinh hoa: Trí thức. Thiếu Nhân hay thiếu Dũng thì Trí chỉ là Trí thức què.

Bởi có NHÂN nên phản ứng với bất nhân, bởi có TRÍ nên phản ứng với sự ngu dốt và loè bịp, bởi có DŨNG nên phản ứng với sự hèn nhát dù là hèn nhát có tính toán, nguỵ trang. Vì thế Trí thức luôn phản biện, luôn gây khó cho một số đối tượng, khiến họ khó chịu, vì Trí thức thường làm rõ mọi điều, xé toạc bức màn che của họ, làm “bể mánh” của họ.

Do đặc trưng ở tính tự do và tính cá thể, trong đó tính thiên bẩm rất quan trọng, nên xã hội có thể dạy kiến thức, dạy cách viết văn… nhưng không thể đào tạo trí thức, đào tạo nhà văn được. Có thể đào tạo một đội ngũ những chuyên viên giỏi, những người có bằng cấp, những người lao động trí óc giỏi, nhưng không thể đào tạo được ra một đội ngũ Trí thức. Trí thức là sự phát triển cá nhân. Việc các Trí thức cùng nhau kết lại để cùng làm một việc gì đó, để đạt một mục tiêu gì đó là việc rất cần thiết nhưng đó chỉ là những việc cụ thể, có giới hạn. Bản chất Trí thức vốn là không đội ngũ, tuy giữa họ luôn có sự hợp đồng, và thành quả của người này tự nhiên kích thích sự sáng tạo ra thành quả của người kia.

Có thể một nghị quyết của Đảng về “đội ngũ Trí thức Xã hội chủ nghĩa” sẽ đưa ra muôn vàn điều ưu ái cho Trí thức, rất đáng quý, nhưng giá được như các nước văn minh, đừng có một nghị quyết về Trí thức thì còn đáng quý hơn, vì đối với những cánh chim thì chẳng cái lồng son nào quý bằng sự tự do. Điều gì cần sự tự nhiên thì càng nhọc công nhào nặn càng làm hỏng việc. Cứ lấy ví dụ trong gia đình, khi đàn con phải họp nhau lại, ra “nghị quyết” từ nay phải thương yêu bố mẹ chẳng hạn (hoặc ngược lại cha mẹ ra “nghị quyết” phải thương yêu con cái) thì liệu bố mẹ có lấy thế làm sung sướng không, hay phải gượng cười để nuốt nước mắt vào trong? Tạo một nền tảng xã hội thuận lợi cho sự phát triển Trí thức là điều rất quý, nhưng đừng xác định Chân lý trước cho họ, đừng định hướng và quản lý cái đầu của họ!

Từ những đặc điểm kể trên hiển nhiên rút ra kết luận: Hiền tài đúng là nguyên khí quốc gia! (Hiền tài đương nhiên là Trí thức).Vốn quý này ai dùng được thì thành công vô hạn, không nhà cầm quyền nào lại không biết điều đó. Chỉ có một điều mà một số người cầm quyền không biết, đó là: Trí thức là NGƯỜI THẦY TỐT nhưng là TÊN ĐẦY TỚ XẤU! (Nếu bị dùng như công cụ, Trí thức chỉ còn là một Trí nô hèn mọn).

Nếu thực tâm coi Trí thức là một người thầy, người thầy ấy sẽ cung cấp không công cho xã hội những hiểu biết quý giá (như số phận con tằm dẫu chẳng ai thuê cũng cứ nhả tơ). Nhưng đã coi họ là thầy của mình thì xin đừng lãnh đạo thầy, đừng chỉ dẫn thầy, đừng khoanh vùng cho thầy, nhất là đừng xoa đầu hay doạ nạt thầy, tội nghiệp! Tồi tệ hơn, nếu dùng Trí thức như một tên đày tớ, mua nó bằng tiền của hay ép nó bằng bạo lực, thì kẻ tay sai này dẫu đắc lực đến mấy, về sau, dù vô tình hay hữu ý nhất định nó sẽ phản thùng, làm cho lão chủ lụn bại không thể chống đỡ, bởi trong trường hợp này Trí thức đã biến tính không còn là Trí thức đúng nghĩa nữa. Khi buộc phải trung thành với chủ, Trí thức không còn trung thành với Chân lý, và sản phẩm của nó chỉ là sản phẩm giả, thay vì phải can gián nó lại toa dập cho chủ vừa lòng.

Trên đây là những ưu điểm đặc trưng của Trí thức với tư cách một tầng lớp tinh hoa. Nhưng một người Trí thức cụ thể thì dễ gì đạt được mức độ lý tưởng ấy. Vả lại cái gì cũng có mặt trái, tập trung mặt này dễ để lại khoảng trống về mặt khác. Ưu điểm này buộc phải chấp nhận nhược điểm khác. Một thiết bị “nano” tinh vi, mỏng mảnh, nhạy cảm thường khó chịu nhiệt, tính chịu quăng quật ắt thua những chất liệu “nồi đồng cối đá”. Đã nhạy cảm với quy luật khách quan thì khó lòng nhắm mắt trung thành.

Những nhược điểm thường thấy ở giới Trí thức cũng là điều cần được nghiên cứu.

*

Ý kiến tôi trong bài này chỉ có thế. Nhưng để thư giãn đôi chút, tôi xin trở lại ý kiến ban đầu về chuyện Công-Nông-Trí: Đảng ta tự xưng là Đảng của giai cấp Công nhân, liên kết với Nông dân là đội quân chủ lực, nên lá cờ Đảng chỉ có liên kết hai thứ búa liềm. Lúc đầu hẳn Đảng phải nghĩ rằng triết lý như thế là hoàn hảo. Này nhá, Công nhân là giai cấp tiền phong tiến bộ nhất, tiêu biểu cho phương thức sản xuất hiện đại, Đảng lại là đội tiền phong tinh hoa của giai cấp Công nhân, lại thêm có lý thuyết Mác-Lê khoa học dẫn đường, thế thì tất cả Trí tuệ đã nằm ở đấy cả rồi, nói Trí thức nữa là thừa. Trí thức chẳng những là khúc “appendice” vô dụng mà đôi khi lại còn bị nhiễm trùng, sưng tấy lên nguy hiểm. Những Trí thức có mặt trong hàng ngũ Cách mạng là Trí thức đã “được đầu hàng giai cấp” (!) tức là đã được thuần hoá, được “công nông hoá” rồi, an toàn rồi.

Nhưng cuộc đời đúng là tên phản biện lại Mác-Lê cực kỳ lì lợm. Mãi rồi cuộc đời cũng lay cho Đảng bừng tỉnh: Ừ, đi vào kỷ nguyên của Trí tuệ mà không trương cờ Trí thức ra thì làm sao thuyết phục thiên hạ? Bởi thế cái anh đầu sỏ cần “đào tận gốc” ngày trước nay đã được phục hoạt, kéo lên cho nhập vào hàng ngũ tiên phong.

Nhưng kẹt nỗi làm sao có thể xếp TRÍ lên trên CÔNG và NÔNG được khi đã “nguyện trung thành” với lý thuyết Cộng sản và lá cờ búa liềm? Vậy là dù có công nhận “Hiền tài là nguyên khí quốc gia”, thì cũng phải để hiền tài, để Trí thức xếp hàng ba sau Công và Nông thôi. Sau khi đã hết lời vinh danh Trí thức thì thứ tự vẫn phải là CÔNG-NÔNG rồi mới đến TRÍ chứ không đảo lộn hàng ngũ được! Bên cạnh đó trong kinh tế thị trường lại còn phải vinh danh Thương nhân bằng cách tổ chức “Ngày Doanh nhân” nữa. (Thế thì cứ chấp nhận luôn cái công thức SĨ-NÔNG-CÔNG THƯƠNG vốn đã có trước đây nhiều thế kỷ cho xong!).

Năm 1996, lần đầu tiên thấy Trí thức được lôi lên đứng sau Công Nông, tôi đã buột miệng làm mấy câu vè:

TRÍ, PHÚ, ĐỊA, HÀO

Bốn anh Trí Phú Địa Hào

Chỉ riêng anh Trí lao đao đến giờ

Đảng ta thương Trí ngu ngơ

Cho CÔNG-NÔNG-TRÍ chung cờ liên minh

Trông lên LIỀM-BÚA hai hình

Trí ta vẫn chẳng thấy mình ở đâu

Quay sang tìm Phú, Địa, Hào

Thấy ba bụng phệ… đã vào… Đảng ta!

HSP-1996

Thế mới biết sửa chữa bằng cách chắp vá là cách chữa cháy khá khôn ngoan, như kiểu “kinh tế thị trường theo định hướng Xã hội chủ nghĩa” là khôn ngoan lắm, nhưng chắp vá Trí thức vào lá cờ Búa Liềm thế nào cho khớp thì đến anh thợ chắp bậc thầy chắc cũng còn nát óc mà nghĩ chưa ra. Bỏ đi, làm cái mới chắc dễ hơn nhiều.

Những ai thực lòng muốn nghe phản biện, thì xin coi đây là một phản biện, liệu có được chăng?

20-7-2008

H. S. P.